\chapter{Удиртгал}
% БАГШИД ҮЗҮҮЛЭХ хуулбар (2026-11-10, ажлын хувилбараас). Засварыг Chapters/Chapter1_Introduction.tex-д хийнэ.
\label{ch:introduction}

%======================================================================
% 2026.11.09-нд хэрэглэгчийн аудитаар дахин бичсэн: систем 1-р давхрын
% ресепшнд байрлана; мэдээллийн эх нь сайт + сургалтын албаны мэдээлэл.
% Бүтэц өмнөх ШМТК-ийн дипломын хэв маягаар (хуулаагүй): үндэслэл ба ач
% холбогдол → зорилго, зорилт → судалгааны асуулт ба шалгуур → хүрээ.
% Шинэ ишлэлүүдийн бүтэн текстийг агент уншсан:
% papers/24-reception-campus-lit-20261109.mn.md
%======================================================================

\section{Үндэслэл, ач холбогдол}
\label{sec:intro-background}

Шинэ Монгол Технологийн Коллежийн нэгдүгээр давхрын ресепшнд өдөр бүр олон хүн мэдээлэл асууж ирдэг. Оюутан, элсэгч, эцэг эх, зочин «энэ өрөө хаана байна», «багш хэзээ хичээлтэй вэ», «элсэлт хэзээ эхлэх вэ», «төлбөр хэд вэ» гэж асууна. Ихэнх асуулт давтагддаг. Хариулт нь сургуулийн вэб сайт, сургалтын албаны хүснэгт, хичээлийн хуваарь гэх мэт олон газарт тархсан байдаг.
2026 оны 11-р сард сургалтын албатай уулзахад ажилтнууд ижил мэдээллийг өөр өөр хүнд дахин дахин тайлбарладаг гэдгийг хэлсэн. Нэгдүгээр давхарт цаасан мэдээлэл өгдөг ч хүмүүс дэлгэрүүлж асуудаг. Мөн сургуулийн вэб сайт, Facebook хуудасны чатбот ба шинэ систем ижил мэдээлэл өгөх ёстой гэсэн шаардлага тавьсан.
% ӨӨРИЙН ҮГЭЭР: уулзалтын огноо, хэнтэй уулзсаныг нэмж зас. Өдөрт хэдэн хүн ирдэг,
% аль асуулт хамгийн их давтагддаг тоо одоогоор алга — авбал энд нэм (агент зохиохгүй).

Ийм асуудал бусад сургуульд ч байдаг. Милано-Бикокка их сургуулийн оюутнууд анги танхим, хичээлийн хуваарь, шалгалтын өрөөний мэдээллийг олоход хэцүү гэж хариулсан \cite{antico2024unimib}. Burch их сургуулийн сайт 470 гаруй хуудастай. Оюутнууд хэрэгтэй мэдээллээ олж чадахгүй үед энгийн асуултаар захиргааны ажилтан руу ханддаг. Энэ нь ажилтны ачааллыг нэмдэг \cite{matic2026burch}. Иймээс олон сургууль байнгын асуултад хариулдаг туслах систем хийж эхэлсэн \cite{trienes2025marcel,tao2025aloha}.

Ресепшнд байрлах туслах дуугаар асуултад хариулах нь зүйтэй. Хүн лангууны өмнө зогсож байхдаа гараараа бичихээс илүү ярьж асуух нь амар. Гэхдээ үүнд гурван бэрхшээл бий.

Нэгдүгээрт, Монгол хэл хэлний технологид нөөц багатай хэл. Олон хэлний Whisper загвар Монгол яриаг 100\%-иас их үгийн алдаатай таньсан \cite{radford2023whisper}. Халх Монгол хэлний нээлттэй сургалтын яриа хэдхэн цаг байдаг \cite{mengke2025khalkha}. Том хэлний загварууд Монгол өгүүлбэрийн дүрмийг бодвол утгыг нь муу ойлгодог \cite{zhang2024mmeval}.

Хоёрдугаарт, олон хүнтэй, чимээтэй газарт яриа таних хэцүү. Карнеги Меллоны их сургуулийн ресепшн робот Valerie-г хийхдээ хонгилын чимээнээс болж яриа танихыг орхиж, гарын товчлуураар бичүүлэхийг сонгосон \cite{gockley2005valerie}. Харин олон нийтийн амралтын өрөөнд туршсан ярианы систем чимээгүй өрөөнийхтэй ойролцоо үр дүн авсан \cite{pappu2014public}.

Гуравдугаарт, хариулт заавал үнэн байх ёстой. Хэлний загвар байхгүй баримт зохиож болно \cite{lewis2020rag}. Сургуулийн туслах буруу холбоос, буруу мэдээлэл өгөхөд оюутнуудын итгэл буурсан \cite{antico2024unimib}. Хэрэглэгчид системийн хариултыг харж шалгахыг хүсдэг \cite{luger2016badpa}. Иймээс хариулт бүр аль эхээс гарсан нь харагдах, эх байхгүй бол «мэдэхгүй» гэж хэлэх хэрэгтэй. Хайлтаар баяжуулсан хариулт (\ac{RAG}) нь эхлээд баримтын сангаас хайж, олсон баримтдаа тулгуурлан хариулдаг \cite{lewis2020rag}. Сургуулийн мэдээлэл өөрчлөгдөхөд загвар сургахгүй, зөвхөн санг шинэчилнэ.

Бидний тоймолсон сургуулийн туслах системүүд их нөөцтэй хэл дээр хийгдсэн бөгөөд хэрэглэгч бичгээр асуудаг \cite{antico2024unimib,trienes2025marcel,tao2025aloha,matic2026burch,gaglio2019campus}. Монгол хэлээр, ресепшнд дуугаар асууж, эх сурвалжтай хариулт авдаг сургуулийн туслах бидний хайлтад олдоогүй (хайлтын аргыг \ref{sec:lit-novelty}-д бичив). Энэ ажлын ач холбогдол үүнд оршино.

\section{Зорилго, зорилт}
\label{sec:intro-goals}

\subsection{Судалгааны зорилго}

Ресепшнд ирсэн хүн сургуулийн мэдээллийг Монгол хэлээр ярьж эсвэл бичиж асуухад сургуулийн албан эхээс хариулт өгч, хариултыг дуугаар уншдаг туслах систем хийж, түүний яриа танилт ба хариултын чанарыг туршилтаар үнэлэх.

\subsection{Судалгааны зорилтууд}

\begin{enumerate}
    \item Ресепшнд ирдэг асуултын төрлийг тодорхойлж, түүнд хэрэгтэй мэдээллийг сургуулийн сайт ба сургалтын албанаас цуглуулж, сэдвээр нь мэдлэгийн модонд бүлэглэх.
    \item Монгол яриаг текст болгох хоёр аргыг ижил бичлэг дээр харьцуулж, системд ашиглах аргыг сонгох.
    \item Асуултад тохирох эхийг хайж, эхийн үгээр хариулдаг, эх олдохгүй бол татгалздаг хариултын хэсгийг хэрэгжүүлэх.
    \item Яриа таних, хайлт, хариулт, дуугаар уншихыг нэг систем болгож, хариултын чанарыг түгжсэн асуултын багцаар хэмжих.
\end{enumerate}

\section{Судалгааны асуулт ба хэмжих шалгуур}
\label{sec:rq}

\begin{enumerate}[label=\textbf{RQ\arabic*.}]
    \item Сургуулийн тухай Монгол асуултыг текст болгоход Монгол хэлээр сургасан Whisper загвар ба Gemini 3.5 Transcribe үйлчилгээ ижил бичлэг дээр ямар алдаатай вэ?
    \item Сургуулийн албан эхээс хайж, дүрмээр эх сонгодог систем эх нь байгаа асуултад хэр зөв хариулж, эх нь байхгүй асуултад хэр татгалзаж байна вэ?
\end{enumerate}

Судалгааны асуулт бүрийг юугаар хэмжихийг \tref{tab:rq_metrics}-т харуулав.

\begin{table}[H]
\centering
\caption{Судалгааны асуулт ба хэмжих шалгуур}
\label{tab:rq_metrics}
\begin{tabular}{|p{1.4cm}|p{6.2cm}|p{6.2cm}|}
\hline
\textbf{Асуулт} & \textbf{Хэмжүүр} & \textbf{Урьдчилж тогтоосон шалгуур} \\
\hline
RQ1 & Үгийн алдааны хувь (\ac{WER}), тэмдэгтийн алдааны хувь (\ac{CER}), нэр, тоо, оны алдаа, хүлээх хугацаа & Gemini-ийн WER Whisper-ээс дор хаяж 10\% бага, CER муудаагүй, Whisper зөв таньсан нэр, тоог алдаагүй байх \\
\hline
RQ2 & Зөв хэсэг хайлтын эхний 1, 3, 5-д орсон тоо; шаардсан утга бүрэн гарсан хариулт; эх нь байхгүй асуултад эхгүйгээр хариулсан тоо & Өөрчлөлт эхгүй хариултыг цөөлж, зөв хариултыг цөөлөхгүй бол л хүлээж авна \\
\hline
\end{tabular}
\end{table}

\section{Хамрах хүрээ}
\label{sec:intro-scope}

Систем сургуулийн албан мэдээллээр л хариулна. Мэдээллийн эх хоёр. Эхнийх нь сургуулийн вэб сайтын нийтэд ил хуудсууд. Хоёр дахь нь сургалтын албанаас авах өрөө, багш, анги, хичээлийн хуваарийн хүснэгт. Хоёр дахь эхийг 2026 оны 11-р сард цуглуулж системд оруулсан. Иймээс эхний хоёр туршилт зөвхөн сайтын мэдээлэл дээр, гурав дахь туршилт хоёр эхийг хамтад нь ашиглаж хийгдсэн. Мөн дуу таних ба хариултын хэсгийг тусад нь туршсан. Дуугаар асуугаад хариулт авах бүтэн урсгалыг хэмжээгүй.

Асуулт ба хариулт Монгол хэлээр байна. Хариултыг шинээр зохиож бичихгүй, сонгосон эхийн үгээр өгнө. Хувийн дүн, ирц зэрэг хувь хүний мэдээлэлд хариулахгүй. Багш өрөөндөө байгаа эсэхийг батлахгүй, зөвхөн нийтэлсэн хуваарийг хэлнэ. Хариултыг дуугаар унших (\ac{TTS}) хэсгийн хүний сонсголын үнэлгээ дараагийн шатанд хийгдэнэ.

\section{Ажлын бүтэц}
\label{sec:intro-structure}

\ref{ch:literature}-р бүлэгт сургуулийн туслах систем, ресепшний туслах, Монгол яриа таних, хайлтаар баяжуулсан хариултын өмнөх судалгааг тоймлов. \ref{ch:method}-р бүлэгт системийн бүтэц, мэдээллийн эх, хоёр туршилтын аргыг бичив. \ref{ch:results}-р бүлэгт туршилтын үр дүн ба тайлбарыг, 5-р бүлэгт дүгнэлт, хязгаарлалт, цаашдын ажлыг бичив.
